\section{Complex Differentiation}

\begin{conceptbox}
    Let $f:\mathbb{C}\to\mathbb{C}$. The derivative of $f$ at $z$ is defined by
    \[
        f'(z)=\lim_{\Delta z\to 0}\frac{f(z+\Delta z)-f(z)}{\Delta z},
    \]
    provided the limit exists and is independent of the direction of approach.
\end{conceptbox}

\subsection*{Motivation / Intuition}

Although the definition resembles that of real calculus, the key difference is that $\Delta z$ can approach $0$ from infinitely many directions in the complex plane.

This makes complex differentiability much more restrictive than real differentiability.

\subsection*{Practical Differentiation Rules}

In practice, once a function is analytic, differentiation proceeds exactly like in real calculus.

\begin{conceptbox}
    If $f$ and $g$ are analytic, then:
    \[
        (f\pm g)'=f'\pm g',
    \]
    \[
        (fg)'=f'g+fg',
    \]
    \[
        \left(\frac{f}{g}\right)'=\frac{f'g-fg'}{g^2}, \quad g\neq 0,
    \]
    \[
        (f\circ g)'=(f'\circ g)\,g'.
    \]
\end{conceptbox}

\subsection*{Analytic Functions}

\begin{conceptbox}
    A function $f$ is \textbf{analytic} at $z_0$ if it is differentiable in a neighborhood of $z_0$.

    If $f$ is analytic everywhere in $\mathbb{C}$, it is called \textbf{entire}.
\end{conceptbox}

\subsection*{Key Observations}

\begin{itemize}
    \item Polynomials are entire.
    \item Rational functions are analytic wherever their denominators are nonzero.
    \item Functions involving $z^*$ are generally not analytic.
\end{itemize}

\begin{examplebox}
    \textbf{Example 1: Differentiate a polynomial}

    Find the derivative of
    \[
        f(z)=z^2+3z+1.
    \]

    Using standard rules:
    \[
        f'(z)=2z+3.
    \]
\end{examplebox}

\begin{examplebox}
    \textbf{Example 2: Differentiate a rational function}

    Find
    \[
        \frac{d}{dz}\left(\frac{1}{z}\right).
    \]

    Rewrite:
    \[
        \frac{1}{z}=z^{-1}.
    \]

    Differentiate:
    \[
        \frac{d}{dz}(z^{-1})=-z^{-2}=-\frac{1}{z^2}, \quad z\neq 0.
    \]
\end{examplebox}

\begin{examplebox}
    \textbf{Example 3: Non-analytic function}

    Let
    \[
        f(z)=|z|^2=x^2+y^2.
    \]

    This function cannot be differentiated using standard complex rules and is not analytic (except trivially at $z=0$).
\end{examplebox}

\begin{noteBox}
    In engineering applications, once analyticity is established, differentiation becomes straightforward and behaves exactly like real calculus. The real challenge is determining whether a function is analytic—this is where the Cauchy--Riemann equations become essential.
\end{noteBox}

\subsection*{Common Mistakes / Notes}

\begin{itemize}
    \item Do not apply differentiation rules unless the function is analytic.
    \item The presence of $z^*$ usually indicates non-analytic behavior.
    \item Complex differentiation is stricter than real differentiation due to direction independence.
\end{itemize}